\documentclass[11pt,oneside]{article}
\usepackage[a4paper,margin=27mm,headheight=15pt]{geometry}
\usepackage[T1]{fontenc}
\usepackage[utf8]{inputenc}
\IfFileExists{libertinus.sty}{\usepackage{libertinus}}{\usepackage{lmodern}}
\usepackage{microtype}
\usepackage{amsmath,amssymb,amsthm,mathtools,mathrsfs}
\usepackage{bm}
\usepackage{booktabs,longtable,array,tabularx}
\usepackage{graphicx}
\usepackage{pgfplots}
\usepackage{xcolor}
\usepackage{enumitem}
\usepackage{fancyhdr}
\usepackage{titlesec}
\usepackage{listings}
\usepackage{xurl}
\usepackage{hyperref}
\usepackage[nameinlink,noabbrev]{cleveref}
\definecolor{linkblue}{RGB}{25,75,135}
\definecolor{shadegray}{RGB}{246,247,249}
\definecolor{rulegray}{RGB}{195,201,208}
\pgfplotsset{compat=1.18}
\hypersetup{
  colorlinks=true,
  linkcolor=linkblue,
  citecolor=linkblue,
  urlcolor=linkblue,
  pdftitle={Further Formulae for the Thue--Morse Sequence},
  pdfsubject={Digital, binomial, valuative, generating-function, finite-difference, Walsh, and product representations},
  pdfkeywords={Thue-Morse, Prouhet, binary digit sum, Lucas theorem, Kummer theorem, Mahler function, Walsh function, Riesz product, Fabius function}
}
\pagestyle{fancy}
\fancyhf{}
\fancyhead[L]{\small Thue--Morse formulae}
\fancyhead[R]{\small\nouppercase{\leftmark}}
\fancyfoot[C]{\thepage}
\renewcommand{\headrulewidth}{0.4pt}
\titleformat{\section}{\Large\bfseries}{\thesection}{0.7em}{}
\titleformat{\subsection}{\large\bfseries}{\thesubsection}{0.7em}{}
\titleformat{\subsubsection}{\normalsize\bfseries}{\thesubsubsection}{0.7em}{}
\setcounter{secnumdepth}{3}
\setcounter{tocdepth}{2}
\setlist[itemize]{leftmargin=2em,itemsep=0.25em,topsep=0.4em}
\setlist[enumerate]{leftmargin=2.3em,itemsep=0.25em,topsep=0.4em}
\setlength{\emergencystretch}{2em}
\newtheorem{theorem}{Theorem}[section]
\newtheorem{proposition}[theorem]{Proposition}
\newtheorem{lemma}[theorem]{Lemma}
\newtheorem{corollary}[theorem]{Corollary}
\newtheorem{identity}[theorem]{Identity}
\theoremstyle{definition}
\newtheorem{definition}[theorem]{Definition}
\newtheorem{algorithm}[theorem]{Algorithm}
\newtheorem{example}[theorem]{Example}
\theoremstyle{remark}
\newtheorem{remark}[theorem]{Remark}
\newtheorem{warning}[theorem]{Warning}
\newcommand{\R}{\mathbb R}
\newcommand{\C}{\mathbb C}
\newcommand{\Q}{\mathbb Q}
\newcommand{\N}{\mathbb N}
\newcommand{\Z}{\mathbb Z}
\newcommand{\Ftwo}{\mathbb F_2}
\newcommand{\e}{\mathrm e}
\newcommand{\ii}{\mathrm i}
\newcommand{\dd}{\,\mathrm d}
\newcommand{\wt}{\operatorname{wt}_2}
\newcommand{\vtwo}{v_2}
\newcommand{\xor}{\mathbin{\oplus}}
\newcommand{\ind}{\mathbf 1}
\newcommand{\coeff}[2]{[z^{#1}]\,#2}
\newcommand{\Bell}{\mathcal B}
\newcommand{\TM}{\mathrm{TM}}
\newcommand{\repo}{\nolinkurl{Analysis/FabiusFunction}}
\newcommand{\lean}[1]{\texttt{#1}}
\newenvironment{proofidea}{\par\smallskip\noindent\textit{Proof architecture.}\ }{\hfill$\square$\par\smallskip}
\newenvironment{boxedremark}{%
  \par\smallskip\noindent\begin{center}\begin{minipage}{0.94\textwidth}%
  \color{black}\hrule\vspace{0.6em}\small
}{%
  \vspace{0.6em}\hrule\end{minipage}\end{center}\smallskip
}
\lstdefinestyle{pseudo}{
  basicstyle=\ttfamily\small,
  backgroundcolor=\color{shadegray},
  frame=single,
  rulecolor=\color{rulegray},
  columns=fullflexible,
  keepspaces=true,
  showstringspaces=false,
  xleftmargin=0.7em,
  xrightmargin=0.7em,
  aboveskip=0.8em,
  belowskip=0.8em
}

\title{\bfseries Further Formulae for the Thue--Morse Sequence\\[0.35em]
\large Digital, binomial, valuative, Mahler-product, finite-difference, Walsh, and infinite-product forms}
\author{Complementary report for the Fabius-function development\\
\small Vladimir Reshetnikov's \texttt{ProveIt} repository}
\date{Repository snapshot consulted: 25 August 2026}

\begin{document}
\maketitle

\begin{abstract}
Section 11.8 of the existing Fabius--Rvachev synthesis gives an unusual
binomial--logarithm formula for the zero--one Thue--Morse bit.  The purpose of
this report is to place that identity inside a much larger, explicitly proved
network of formulae.  We derive representations of the signed value
\(\varepsilon_n=(-1)^{\wt(n)}\) and the bit
\(\tau(n)=\wt(n)\bmod 2\) from binary blocks, floor sums, factorial
valuations, central binomial coefficients, sparse entries of Pascal's
triangle, coefficient extraction, finite differences, automata, tensor
products, Walsh characters, discrete Fourier transforms, Riesz products, and
self-similar infinite products.

Several identities appear to be especially useful as additions to the formal
library.  These include a carry-corrected addition law
\[
 \varepsilon_{a+b}=\varepsilon_a\varepsilon_b
 (-1)^{\vtwo\binom{a+b}{a}},
\]
a power-of-two Pascal product
\[
 \varepsilon_n=(-1)^{\sum_j\binom{n}{2^j}}
 =\prod_j\cos\!\left(\pi\binom{n}{2^j}\right),
\]
an exact formula for every ordinary prefix sum, and an all-orders closed
formula for the Prouhet power sums
\(\sum_{n<2^m}\varepsilon_n n^r\).  The latter is given both as a finite sum
over compositions of \(r\) and as a complete Bell-polynomial expression.
We also obtain exact Fourier formulae and a compact product formula for every
base-\(b\) Thue--Morse constant.

All proofs are supplied in ordinary mathematics.  The phrase ``derived here''
means derived in this report from standard identities; it is not a claim of
priority in the literature.  The final section identifies natural theorem
statements and proof dependencies for a future Lean formalization alongside
the existing modules \texttt{DyadicClosedForm.lean},
\texttt{ThueMorsePrefix.lean}, and the primitive-recursive bit evaluator.
\end{abstract}

\tableofcontents
\newpage

\section{Scope, conventions, and the connection with the Fabius function}
\label{sec:scope}

The signed Thue--Morse sequence enters the Fabius-function theory through the
global translate expansion and through the exact cancellation of dyadic
blocks.  In the notation of the main synthesis \cite{proveitSynthesis},
\begin{equation}
 \varepsilon_n=(-1)^{\wt(n)},
 \qquad
 \tau(n)=\wt(n)\bmod 2,
 \qquad
 \varepsilon_n=1-2\tau(n).
 \label{eq:basic-def}
\end{equation}
Here \(\wt(n)\) is the number of ones in the binary expansion of \(n\).  The
existing development proves the binary recurrences, block concatenation,
Prouhet cancellation, exact iterated-prefix identities, and the
binomial--logarithm formula quoted in the question
\cite{proveit,proveitSynthesis}.  We retain those results as the base layer and
look for alternative formulae that expose different structures.

Write
\begin{equation}
 n=\sum_{j\ge0}b_j(n)2^j,
 \qquad b_j(n)\in\{0,1\},
 \label{eq:binary-digits}
\end{equation}
where all but finitely many digits vanish.  We use \(\vtwo(M)\) for the
exponent of \(2\) in a positive integer \(M\), and set \(\vtwo(0!)=0\).
For bitwise exclusive-or we write \(a\xor b\).  Empty products are one.

\begin{boxedremark}
\textbf{Two forms, one parity character.}
The zero--one sequence is natural for substitutions and constants whose
base expansion consists of Thue--Morse digits.  The signed sequence is natural
for products, finite differences, Fourier analysis, and the Fabius translate
sum.  Every formula below can be moved between the two by
\(\tau(n)=(1-\varepsilon_n)/2\).
\end{boxedremark}

The report is organized around four ``engines'' that repeatedly generate new
identities.
\begin{enumerate}
\item Binary digits turn \(\varepsilon_n\) into a character of a Boolean cube.
\item Legendre--Kummer valuation formulae translate digit sums into factorial
      and binomial arithmetic.
\item The finite polynomial
      \(\sum_{n<2^m}\varepsilon_nz^n\) factors completely.
\item That factorization is simultaneously a generating-function identity, a
      mixed finite difference, and a Fourier product.
\end{enumerate}

\section{A master catalogue of individual-term formulae}
\label{sec:catalogue}

The next theorem gathers formulae that compute a single term.  Later sections
separate them and explain what each representation is good for.

\begin{theorem}[Equivalent formulae for one Thue--Morse sign]
\label{thm:master-catalogue}
Let \(n\ge0\), and choose \(m\) with \(n<2^m\).  Let
\begin{equation}
 \nu(n)=\#\left\{0\le k\le n:\binom nk\text{ is odd}\right\}.
 \label{eq:nu-row}
\end{equation}
Then all of the following quantities equal \(\varepsilon_n\):
\begin{align}
 &(-1)^{\sum_{j=0}^{m-1}b_j(n)},
 &&\prod_{j=0}^{m-1}\bigl(1-2b_j(n)\bigr),
 \label{eq:catalog-digits}\\
 &(-1)^{\,n-\sum_{j\ge1}\lfloor n/2^j\rfloor},
 &&(-1)^{\,n-\vtwo(n!)},
 \label{eq:catalog-floor-factorial}\\
 &(-1)^{\vtwo\binom{2n}{n}},
 &&(-1)^{\sum_{j=0}^{m-1}\binom{n}{2^j}},
 \label{eq:catalog-binomial}\\
 &\prod_{j=0}^{m-1}\cos\!\left(\pi\binom{n}{2^j}\right),
 &&(-1)^{\log_2\nu(n)},
 \label{eq:catalog-pascal}\\
 &\prod_{j\ge0}\cos\!\left(\pi\left\lfloor\frac{n}{2^j}\right\rfloor\right),
 &&[z^n]\prod_{j\ge0}\bigl(1-z^{2^j}\bigr).
 \label{eq:catalog-cos-coeff}
\end{align}
Every infinite product in \eqref{eq:catalog-cos-coeff} is effectively finite:
its factors are one, or irrelevant to the coefficient, beyond the binary
length of \(n\).
\end{theorem}

\begin{proof}
The first pair is the definition of the parity character of the binary digit
vector.  Legendre's formula gives
\begin{equation}
 \vtwo(n!)=\sum_{j\ge1}\left\lfloor\frac{n}{2^j}\right\rfloor
          =n-\wt(n),
 \label{eq:legendre-base-two}
\end{equation}
which proves \eqref{eq:catalog-floor-factorial}.  Applying the same formula to
\((2n)!/(n!)^2\) yields
\begin{equation}
 \vtwo\binom{2n}{n}
 =2\wt(n)-\wt(2n)=\wt(n),
 \label{eq:central-valuation}
\end{equation}
so the first formula in \eqref{eq:catalog-binomial} follows.

Lucas's theorem modulo two says
\begin{equation}
 \binom{n}{2^j}\equiv b_j(n)\pmod2.
 \label{eq:lucas-power-two}
\end{equation}
Therefore the parity of the second exponent in
\eqref{eq:catalog-binomial} is \(\sum_jb_j(n)=\wt(n)\).  Replacing
\((-1)^A\) by \(\cos(\pi A)\) gives the first formula in
\eqref{eq:catalog-pascal}.  Glaisher's count, itself an immediate consequence
of Lucas's theorem, is \(\nu(n)=2^{\wt(n)}\); hence
\(\log_2\nu(n)=\wt(n)\).

For the floor-cosine product, observe that
\begin{equation}
 \sum_{j\ge0}\left\lfloor\frac{n}{2^j}\right\rfloor
 =n+\vtwo(n!)=2n-\wt(n).
 \label{eq:all-floor-sum}
\end{equation}
Its parity is therefore the parity of \(\wt(n)\).  Finally, expanding
\(\prod_j(1-z^{2^j})\) amounts to choosing a subset of powers of two.  Binary
expansion makes that subset unique, and its coefficient is
\((-1)^{\#\text{chosen powers}}=(-1)^{\wt(n)}\).
\end{proof}

\begin{corollary}[Trigonometric zero--one formulae]
\label{cor:trig-bit}
For every \(n\ge0\),
\begin{align}
 \tau(n)
 &=\sin^2\!\left(\frac{\pi}{2}\bigl(n-\vtwo(n!)\bigr)\right)
 \label{eq:tau-factorial-sine}\\
 &=\sin^2\!\left(\frac{\pi}{2}\,
          \vtwo\binom{2n}{n}\right)
 \label{eq:tau-central-sine}\\
 &=\frac12\left(1-
      \prod_{j\ge0}\cos\!\left(\pi\binom{n}{2^j}\right)\right).
 \label{eq:tau-pascal-cos}
\end{align}
\end{corollary}

\begin{proof}
For an integer \(A\), \(\sin^2(\pi A/2)\) is zero when \(A\) is even and one
when \(A\) is odd.  Apply \cref{thm:master-catalogue} and
\(\tau=(1-\varepsilon)/2\).
\end{proof}

\begin{remark}
Formula \eqref{eq:tau-pascal-cos} is a direct companion to the
binomial--logarithm formula in the existing article.  Instead of inspecting
all \(n+1\) entries of one row and then taking a logarithm, it inspects only
the \(O(\log n)\) entries whose column indices are powers of two.  Lucas's
theorem makes those sparse entries exactly the binary digits of the row
index.
\end{remark}

\section{Digital blocks, complements, exclusive-or, and carries}
\label{sec:digital}

\subsection{Recurrences and block concatenation}

The familiar two-term recursion is
\begin{equation}
 \varepsilon_{2q}=\varepsilon_q,
 \qquad
 \varepsilon_{2q+1}=-\varepsilon_q,
 \label{eq:evenodd-sign}
\end{equation}
with the zero--one version
\begin{equation}
 \tau(2q)=\tau(q),
 \qquad
 \tau(2q+1)=1-\tau(q).
 \label{eq:evenodd-bit}
\end{equation}
Iterating it gives a stronger formula that is often the right replacement for
case splitting on one final bit.

\begin{theorem}[Digit-block multiplicativity]
\label{thm:block-multiplicativity}
For \(q,m\ge0\) and \(0\le r<2^m\),
\begin{equation}
 \varepsilon_{q2^m+r}=\varepsilon_q\varepsilon_r,
 \qquad
 \tau(q2^m+r)=\tau(q)\xor\tau(r).
 \label{eq:block-multiplicativity}
\end{equation}
\end{theorem}

\begin{proof}
The binary digits of \(q2^m\) and \(r\) occupy disjoint positions, so
\(\wt(q2^m+r)=\wt(q)+\wt(r)\).  Taking parity proves both identities.
\end{proof}

This theorem simultaneously describes every member of the \(2\)-kernel.  For
a fixed residue \(r<2^m\), the subsequence along that residue class is either
the original signed sequence or its negative:
\begin{equation}
 \bigl(\varepsilon_{q2^m+r}\bigr)_{q\ge0}
 =\varepsilon_r\bigl(\varepsilon_q\bigr)_{q\ge0}.
 \label{eq:kernel-residue}
\end{equation}
For zero--one values it is either \(\tau\) or its bitwise complement.

\subsection{Reflection inside a dyadic block}

\begin{proposition}[Dyadic reflection]
\label{prop:dyadic-reflection}
If \(0\le r<2^m\), then
\begin{equation}
 \varepsilon_{2^m-1-r}=(-1)^m\varepsilon_r,
 \qquad
 \tau(2^m-1-r)=\tau(r)\xor(m\bmod2).
 \label{eq:dyadic-reflection}
\end{equation}
\end{proposition}

\begin{proof}
Within exactly \(m\) binary positions, \(2^m-1-r\) is the bitwise complement
of \(r\).  Hence its digit sum is \(m-\wt(r)\).
\end{proof}

Thus the word of length \(2^m\) is a palindrome when \(m\) is even and an
anti-palindrome when \(m\) is odd.  The adjacent dyadic block is always the
negative of the first:
\begin{equation}
 \varepsilon_{2^m+r}=-\varepsilon_r
 \qquad(0\le r<2^m).
 \label{eq:next-block-negative}
\end{equation}

\subsection{Exclusive-or is the natural group law}

\begin{theorem}[Walsh-character law]
\label{thm:xor-law}
For all nonnegative integers \(a,b\),
\begin{equation}
 \varepsilon_{a\xor b}=\varepsilon_a\varepsilon_b,
 \qquad
 \tau(a\xor b)=\tau(a)\xor\tau(b).
 \label{eq:xor-law}
\end{equation}
\end{theorem}

\begin{proof}
At each bit, exclusive-or is addition in \(\Ftwo\).  Therefore
\(\wt(a\xor b)\equiv\wt(a)+\wt(b)\pmod2\).
\end{proof}

This says that \(\varepsilon\) is a one-dimensional character of the infinite
Boolean group of finite binary strings.  Ordinary addition differs from
exclusive-or precisely by carries, and Kummer's theorem measures those
carries.

\subsection{A carry-corrected addition law}

\begin{theorem}[Addition formula through a binomial valuation]
\label{thm:addition-law}
For all \(a,b\ge0\),
\begin{equation}
 \boxed{\quad
 \varepsilon_{a+b}=\varepsilon_a\varepsilon_b
          (-1)^{\vtwo\binom{a+b}{a}}.
 \quad}
 \label{eq:addition-law}
\end{equation}
Equivalently,
\begin{equation}
 \tau(a+b)=\tau(a)\xor\tau(b)\xor
       \left(\vtwo\binom{a+b}{a}\bmod2\right).
 \label{eq:addition-law-bit}
\end{equation}
\end{theorem}

\begin{proof}
Legendre's formula gives the base-two form of Kummer's identity,
\begin{equation}
 \vtwo\binom{a+b}{a}
 =\wt(a)+\wt(b)-\wt(a+b).
 \label{eq:kummer-digit}
\end{equation}
Rearranging this equality modulo two and exponentiating \((-1)\) proves
\eqref{eq:addition-law}.
\end{proof}

Kummer's theorem identifies the valuation in \eqref{eq:kummer-digit} with the
number of binary carries, counted with propagation.  Consequently
\(\varepsilon\) is multiplicative under ordinary addition exactly when that
carry count is even.  In particular, carry-free addition gives
\begin{equation}
 (a\mathbin{\&}b=0)
 \quad\Longrightarrow\quad
 \varepsilon_{a+b}=\varepsilon_a\varepsilon_b.
 \label{eq:carry-free}
\end{equation}
The block formula \eqref{eq:block-multiplicativity} is the most important
special case.

\begin{corollary}[Multinomial addition law]
\label{cor:multinomial-addition}
Let \(a_1,\ldots,a_s\ge0\), and put \(N=a_1+\cdots+a_s\).  Then
\begin{equation}
 \varepsilon_N
 =\left(\prod_{j=1}^s\varepsilon_{a_j}\right)
  (-1)^{\vtwo\binom{N}{a_1,\ldots,a_s}},
 \label{eq:multinomial-addition}
\end{equation}
where \(\binom{N}{a_1,\ldots,a_s}=N!/(a_1!\cdots a_s!)\).
\end{corollary}

\begin{proof}
Legendre's formula yields
\(
 \vtwo\binom{N}{a_1,\ldots,a_s}
 =\sum_j\wt(a_j)-\wt(N)
\), and the parity argument is unchanged.
\end{proof}

\begin{example}
Take \(a=13=(1101)_2\) and \(b=7=(0111)_2\).  Their signs are
\(\varepsilon_{13}=-1\) and \(\varepsilon_7=-1\), while
\(a+b=20=(10100)_2\) has sign \(+1\).  Formula
\eqref{eq:addition-law} predicts an even carry count.  Indeed
\begin{equation*}
 \vtwo\binom{20}{13}
 =\wt(13)+\wt(7)-\wt(20)=3+3-2=4.
\end{equation*}
\end{example}

\section{Pascal-triangle formulae beyond the binomial--logarithm identity}
\label{sec:pascal}

\subsection{Power-of-two columns recover the binary digits}

Lucas's congruence modulo two has a particularly sharp specialization.  The
integer \(2^j\) has only one nonzero binary digit, so
\begin{equation}
 \binom{n}{2^j}\equiv b_j(n)\pmod2.
 \label{eq:power-two-column}
\end{equation}
Thus the parity pattern of the sparse columns
\(1,2,4,8,\ldots\) in the \(n\)th row is literally the binary expansion of
\(n\).

\begin{theorem}[Sparse Pascal-row formula]
\label{thm:sparse-pascal}
For every \(n\ge0\),
\begin{align}
 \varepsilon_n
 &=(-1)^{\sum_{j\ge0}\binom{n}{2^j}}
 \label{eq:sparse-pascal-exp}\\
 &=\prod_{j\ge0}(-1)^{\binom{n}{2^j}}
 =\prod_{j\ge0}\cos\!\left(\pi\binom{n}{2^j}\right),
 \label{eq:sparse-pascal-product}
\end{align}
and
\begin{equation}
 \tau(n)=\frac{1-\prod_{j\ge0}
          \cos\!\left(\pi\binom{n}{2^j}\right)}{2}.
 \label{eq:sparse-pascal-bit}
\end{equation}
\end{theorem}

\begin{proof}
Only the parity of each exponent matters.  Apply
\eqref{eq:power-two-column} and sum over the binary positions.
\end{proof}

This is computationally different from the formula in Section 11.8 of the
main document.  The logarithmic formula scans all entries of the row, but
needs no distinguished columns.  The sparse formula scans only logarithmically
many entries, but singles out the columns \(2^j\).

\subsection{The Glaisher count in compressed form}

Let \(\nu(n)\) be the number of odd entries in row \(n\).  Lucas's theorem
shows that \(k\) is allowed at each one-bit of \(n\) and forced to zero at each
zero-bit; hence
\begin{equation}
 \nu(n)=2^{\wt(n)}.
 \label{eq:glaisher-again}
\end{equation}
The most direct logarithmic formula is therefore
\begin{equation}
 \varepsilon_n=(-1)^{\log_2\nu(n)},
 \qquad
 \tau(n)=\frac{1-(-1)^{\log_2\nu(n)}}{2}.
 \label{eq:direct-glaisher-log}
\end{equation}
The existing document encodes \(\nu(n)\) without using an oddness indicator.
If
\begin{equation}
 \mathcal X(n)=\sum_{k=0}^n(-1)^{\binom nk},
 \label{eq:X-repeat}
\end{equation}
then each even entry contributes \(+1\) and each odd entry contributes \(-1\),
so
\begin{equation}
 \nu(n)=\frac{n+1-\mathcal X(n)}{2}.
 \label{eq:nu-from-X}
\end{equation}
Substituting this into \eqref{eq:direct-glaisher-log} gives exactly the
binomial--logarithm formula, including its one-unit shift in the exponent.
This decomposition makes clear that the essential invariant is the row's odd
entry count; \(\mathcal X(n)\) is one elegant way to recover that count.

\subsection{Central binomial coefficients and doubling carries}

\begin{proposition}[Central-binomial valuation formula]
\label{prop:central-binomial}
For every \(n\ge0\),
\begin{equation}
 \vtwo\binom{2n}{n}=\wt(n),
 \qquad
 \varepsilon_n=(-1)^{\vtwo\binom{2n}{n}},
 \qquad
 \tau(n)=\vtwo\binom{2n}{n}\bmod2.
 \label{eq:central-binomial-formula}
\end{equation}
\end{proposition}

\begin{proof}
Either use \eqref{eq:central-valuation}, or apply Kummer's theorem to the
addition \(n+n\).  The number of carries in that addition, with propagation,
is \(\wt(n)\).
\end{proof}

The formula is useful when a proof already contains central binomial
coefficients or Catalan-type expressions.  It also gives a term formula with
no explicit binary digits, floor functions, or row-wide sum.

\section{Floor sums, factorial valuations, and the ruler function}
\label{sec:valuations}

\subsection{Legendre's formula as a digit extractor}

From \eqref{eq:legendre-base-two},
\begin{equation}
 \wt(n)=n-\vtwo(n!)
       =n-\sum_{j\ge1}\left\lfloor\frac{n}{2^j}\right\rfloor.
 \label{eq:weight-floor}
\end{equation}
Consequently
\begin{equation}
 \varepsilon_n
 =(-1)^{n-\vtwo(n!)}
 =(-1)^{n+\vtwo(n!)}.
 \label{eq:factorial-sign}
\end{equation}
The last equality uses only the fact that subtraction and addition coincide
modulo two.  A floor-only form is
\begin{equation}
 \varepsilon_n
 =(-1)^{\,n-\sum_{j\ge1}\lfloor n/2^j\rfloor}.
 \label{eq:floor-only-sign}
\end{equation}

The version with all floors beginning at \(j=0\) is unexpectedly symmetric:
\begin{equation}
 \varepsilon_n
 =(-1)^{\sum_{j\ge0}\lfloor n/2^j\rfloor}
 =\prod_{j\ge0}\cos\!\left(\pi\left\lfloor\frac{n}{2^j}\right\rfloor\right).
 \label{eq:floor-cos-product}
\end{equation}
Indeed the exponent is \(2n-\wt(n)\), which has the same parity as
\(\wt(n)\).

\subsection{A local recurrence controlled by trailing zeros}

Let \(\rho(k)=\vtwo(k)\) for \(k\ge1\); this is the classical ruler function.
Incrementing \(n\) to \(n+1\) changes a terminal string of ones into zeros and
changes the preceding zero into a one.  Therefore
\begin{equation}
 \wt(n+1)-\wt(n)=1-\vtwo(n+1).
 \label{eq:weight-increment}
\end{equation}

\begin{theorem}[Ruler-function update]
\label{thm:ruler-update}
For every \(n\ge0\),
\begin{equation}
 \varepsilon_{n+1}
 =-(-1)^{\vtwo(n+1)}\varepsilon_n.
 \label{eq:ruler-update-sign}
\end{equation}
Equivalently, the zero--one value flips exactly when \(\vtwo(n+1)\) is even:
\begin{equation}
 \tau(n+1)=\tau(n)\xor
       \bigl(1+\vtwo(n+1)\bmod2\bigr).
 \label{eq:ruler-update-bit}
\end{equation}
\end{theorem}

\begin{proof}
Raise \((-1)\) to both sides of \eqref{eq:weight-increment}.
\end{proof}

Telescoping \eqref{eq:ruler-update-sign} from \(0\) to \(n\) gives another
factorial formula:
\begin{equation}
 \varepsilon_n
 =\prod_{k=1}^n(-1)^{1+\vtwo(k)}
 =(-1)^{n+\sum_{k=1}^n\vtwo(k)}
 =(-1)^{n+\vtwo(n!)}.
 \label{eq:ruler-product}
\end{equation}
This formula generates the sequence sequentially without extracting any
binary digits of the current index.  It is attractive in streaming
algorithms, where \(\vtwo(n+1)\) can be read from the least significant set
bit.

\section{Finite and infinite generating products}
\label{sec:generating}

\subsection{The finite Prouhet polynomial}

For \(m\ge0\), define
\begin{equation}
 P_m(z)=\sum_{n=0}^{2^m-1}\varepsilon_nz^n.
 \label{eq:P-def}
\end{equation}

\begin{theorem}[Complete factorization]
\label{thm:finite-product}
For every \(m\ge0\),
\begin{equation}
 P_m(z)=\prod_{j=0}^{m-1}\bigl(1-z^{2^j}\bigr).
 \label{eq:finite-product}
\end{equation}
\end{theorem}

\begin{proof}
The subset expansion of the product is
\begin{equation*}
 \sum_{S\subseteq\{0,\ldots,m-1\}}(-1)^{|S|}
 z^{\sum_{j\in S}2^j}.
\end{equation*}
Every integer below \(2^m\) has a unique such subset, and \(|S|=\wt(n)\).
\end{proof}

The factorization also follows from the word recursion
\(P_{m+1}(z)=P_m(z)-z^{2^m}P_m(z)\).  Combining it with block
multiplicativity gives a shifted block formula:
\begin{equation}
 \sum_{r=0}^{2^m-1}\varepsilon_{q2^m+r}z^r
 =\varepsilon_qP_m(z).
 \label{eq:block-polynomial}
\end{equation}

\subsection{The Mahler product and functional equation}

For \(|z|<1\), or coefficientwise as a formal power series, let
\begin{equation}
 E(z)=\sum_{n\ge0}\varepsilon_nz^n.
 \label{eq:E-def}
\end{equation}
Taking the coefficientwise limit of \eqref{eq:finite-product} gives
\begin{equation}
 \boxed{\quad E(z)=\prod_{j\ge0}\bigl(1-z^{2^j}\bigr).\quad}
 \label{eq:infinite-product-E}
\end{equation}
Peeling off the first factor yields the first-order Mahler equation
\begin{equation}
 E(z)=(1-z)E(z^2).
 \label{eq:mahler-E}
\end{equation}
More generally, the \(2^m\)-dissection is
\begin{equation}
 E(z)=P_m(z)E(z^{2^m}).
 \label{eq:dissection-E}
\end{equation}

For the zero--one generating function
\begin{equation}
 T(z)=\sum_{n\ge0}\tau(n)z^n,
 \label{eq:T-def}
\end{equation}
we have
\begin{equation}
 T(z)=\frac{1}{2(1-z)}-\frac12E(z),
 \label{eq:T-from-E}
\end{equation}
and therefore
\begin{equation}
 T(z)=(1-z)T(z^2)+\frac{z}{1-z^2}.
 \label{eq:mahler-T}
\end{equation}
Mahler used this self-similarity to prove transcendence results for values at
algebraic points inside the unit disk \cite{mahler1929}.

\subsection{Every base-\texorpdfstring{\(b\)}{b} Thue--Morse constant}

For a real \(b>1\), define the base-\(b\) digit constant
\begin{equation}
 \xi_b=\sum_{n\ge0}\frac{\tau(n)}{b^{n+1}}.
 \label{eq:xi-b-def}
\end{equation}
This is the real number whose base-\(b\) fractional digits are the Thue--Morse
bits when \(b\) is an integer.

\begin{theorem}[Product formula for \(\xi_b\)]
\label{thm:base-b-constant}
For every real \(b>1\),
\begin{equation}
 \boxed{\quad
 \xi_b=\frac{1}{2(b-1)}-
       \frac{1}{2b}\prod_{j\ge0}\left(1-b^{-2^j}\right).
 \quad}
 \label{eq:base-b-constant}
\end{equation}
\end{theorem}

\begin{proof}
Set \(z=1/b\) in \eqref{eq:T-from-E} and multiply by \(z\).
\end{proof}

For example,
\begin{align*}
 \xi_2&=0.412454033640107597783361368258455\ldots,\\
 \xi_3&=0.152468764055948430438283372398260\ldots,\\
 \xi_{10}&=0.0110100110010110069261571648290206\ldots.
\end{align*}
The first number is the classical Prouhet--Thue--Morse constant.  Formula
\eqref{eq:base-b-constant} makes the signed product and the zero--one digit
expansion two evaluations of the same Mahler function.

\section{Exact prefix sums and optimal balance bounds}
\label{sec:prefix}

The first iterated prefix sum in the existing development vanishes at every
odd endpoint.  With an exclusive endpoint, there is an especially compact
formula valid at every integer.

\begin{definition}[Exclusive prefix sum]
\label{def:A-prefix}
Set
\begin{equation}
 A(N)=\sum_{n=0}^{N-1}\varepsilon_n
 \qquad(N\ge0).
 \label{eq:A-def}
\end{equation}
\end{definition}

\begin{theorem}[Exact prefix formula]
\label{thm:exact-prefix}
For every \(m\ge0\),
\begin{equation}
 A(2m)=0,
 \qquad
 A(2m+1)=\varepsilon_m.
 \label{eq:prefix-cases}
\end{equation}
Equivalently,
\begin{equation}
 \boxed{\quad
 A(N)=\frac{1-(-1)^N}{2}\,\varepsilon_{\lfloor N/2\rfloor}.
 \quad}
 \label{eq:prefix-compact}
\end{equation}
\end{theorem}

\begin{proof}
Pair adjacent terms using
\(\varepsilon_{2q}+\varepsilon_{2q+1}=0\).  An odd-length prefix has one
unpaired final term \(\varepsilon_{2m}=\varepsilon_m\).
\end{proof}

The inclusive sum through \(N\) is therefore
\begin{equation}
 \sum_{n=0}^{N}\varepsilon_n
 =\frac{1+(-1)^N}{2}\,\varepsilon_{\lfloor N/2\rfloor}.
 \label{eq:inclusive-prefix}
\end{equation}

\begin{corollary}[Exact number of ones in a prefix]
\label{cor:prefix-ones}
Let \(C(N)=\sum_{n=0}^{N-1}\tau(n)\).  Then
\begin{equation}
 C(N)=\frac{N-A(N)}{2},
 \label{eq:C-from-A}
\end{equation}
and in cases
\begin{equation}
 C(2m)=m,
 \qquad
 C(2m+1)=m+\tau(m).
 \label{eq:prefix-one-cases}
\end{equation}
\end{corollary}

Thus every even-length prefix is exactly balanced.  Every odd-length prefix
has imbalance exactly one.

\begin{corollary}[Uniform interval discrepancy]
\label{cor:interval-discrepancy}
For all integers \(0\le a\le b\),
\begin{equation}
 \left|\sum_{n=a}^{b-1}\varepsilon_n\right|\le2,
 \qquad
 \left|\sum_{n=a}^{b-1}\tau(n)-\frac{b-a}{2}\right|\le1.
 \label{eq:interval-discrepancy}
\end{equation}
Both constants are attained.
\end{corollary}

\begin{proof}
The signed interval sum is \(A(b)-A(a)\), and
\(A(N)\in\{-1,0,1\}\).  The zero--one estimate follows from
\(\tau=(1-\varepsilon)/2\).  For example,
\(\varepsilon_5+\varepsilon_6=2\), so the signed bound is sharp.
\end{proof}

Aligned dyadic blocks are more rigid.  If \(m\ge1\), then
\begin{equation}
 \sum_{r=0}^{2^m-1}\varepsilon_{q2^m+r}
 =\varepsilon_q\sum_{r=0}^{2^m-1}\varepsilon_r=0.
 \label{eq:aligned-block-zero}
\end{equation}

\section{The mixed finite-difference master identity}
\label{sec:finite-difference}

The Prouhet cancellation is usually stated for monomials.  The factorization
\eqref{eq:finite-product} gives a stronger identity for an arbitrary test
function.

For a function \(f\), let
\begin{equation}
 (\Delta_hf)(x)=f(x+h)-f(x).
 \label{eq:delta-def}
\end{equation}

\begin{theorem}[Boolean-cube finite difference]
\label{thm:finite-difference-master}
For every \(m\ge0\) and every function for which the displayed values are
defined,
\begin{equation}
 \boxed{\quad
 \sum_{n=0}^{2^m-1}\varepsilon_nf(x+n)
 =(-1)^m\bigl(\Delta_1\Delta_2\Delta_4\cdots
       \Delta_{2^{m-1}}f\bigr)(x).
 \quad}
 \label{eq:finite-difference-master}
\end{equation}
\end{theorem}

\begin{proof}
Let \(E_hf(x)=f(x+h)\).  Expanding
\(\prod_{j=0}^{m-1}(I-E_{2^j})\) gives one term for every subset of the binary
steps.  The subset sum is the corresponding integer \(n<2^m\), and its sign
is \((-1)^{\wt(n)}\).  Since \(I-E_h=-\Delta_h\), the right side acquires the
factor \((-1)^m\).
\end{proof}

\subsection{An exact integral cubature formula}

If \(f\) is \(m\) times continuously differentiable on the relevant interval,
iterated use of the fundamental theorem of calculus gives
\begin{equation}
 \Delta_{h_1}\cdots\Delta_{h_m}f(x)
 =h_1\cdots h_m
  \int_{[0,1]^m}f^{(m)}\!\left(x+\sum_{j=1}^mh_ju_j\right)
  \dd\bm u.
 \label{eq:iterated-difference-integral}
\end{equation}
Taking \(h_j=2^{j-1}\) yields the following exact identity.

\begin{corollary}[Integral representation]
\label{cor:integral-difference}
If \(f\in C^m\), then
\begin{equation}
 \sum_{n=0}^{2^m-1}\varepsilon_nf(x+n)
 =(-1)^m2^{\binom m2}
  \int_{[0,1]^m}
  f^{(m)}\!\left(x+\sum_{j=0}^{m-1}2^ju_j\right)\dd\bm u.
 \label{eq:integral-difference}
\end{equation}
\end{corollary}

This formula explains both the cancellation order and the sharp surviving
coefficient without expanding any powers.

\begin{corollary}[Prouhet extraction]
\label{cor:prouhet-extraction}
If \(p\) is a polynomial of degree below \(m\), then
\begin{equation}
 \sum_{n<2^m}\varepsilon_np(x+n)=0.
 \label{eq:prouhet-general}
\end{equation}
If \(p(y)=cy^m+\text{lower-degree terms}\), then
\begin{equation}
 \sum_{n<2^m}\varepsilon_np(x+n)
 =(-1)^m2^{\binom m2}m!\,c.
 \label{eq:prouhet-leading}
\end{equation}
\end{corollary}

The translation \(x\) disappears at the sharp degree.  Taking
\(p(y)=y^r\) recovers the Prouhet identities already used by the Fabius
formalization.

\section{All higher dyadic power sums}
\label{sec:power-sums}

The vanishing range and the first nonzero power are only the beginning.  The
same product gives a finite closed formula for every higher power.

\begin{definition}[Dyadic Thue--Morse power sum]
For \(m,r\ge0\), set
\begin{equation}
 S_{m,r}=\sum_{n=0}^{2^m-1}\varepsilon_n n^r.
 \label{eq:S-def}
\end{equation}
We use the convention \(0^0=1\).
\end{definition}

\subsection{A composition formula}

Substitute \(z=\e^x\) in \eqref{eq:finite-product}:
\begin{equation}
 \sum_{r\ge0}S_{m,r}\frac{x^r}{r!}
 =\sum_{n<2^m}\varepsilon_n\e^{nx}
 =\prod_{j=0}^{m-1}\bigl(1-\e^{2^jx}\bigr).
 \label{eq:moment-egf}
\end{equation}
Expanding each factor into its exponential series gives the promised
all-orders expression.

\begin{theorem}[All-order composition formula]
\label{thm:power-sum-composition}
Let \(m\ge1\).  For every \(r\ge0\),
\begin{equation}
 \boxed{\quad
 S_{m,r}=(-1)^m r!
 \sum_{\substack{q_0,\ldots,q_{m-1}\ge1\\
                  q_0+\cdots+q_{m-1}=r}}
 \frac{2^{\sum_{j=0}^{m-1}jq_j}}
      {q_0!q_1!\cdots q_{m-1}!}.
 \quad}
 \label{eq:power-sum-composition}
\end{equation}
When \(r<m\), the sum is empty and equals zero.
\end{theorem}

\begin{proof}
Write
\begin{equation*}
 1-\e^{2^jx}=-\sum_{q_j\ge1}\frac{2^{jq_j}x^{q_j}}{q_j!}
\end{equation*}
and extract the coefficient of \(x^r\) from the product in
\eqref{eq:moment-egf}.
\end{proof}

At \(r=m\), every \(q_j\) must equal one, and
\begin{equation}
 S_{m,m}=(-1)^m m!2^{0+1+\cdots+(m-1)}
 =(-1)^m2^{\binom m2}m!,
 \label{eq:power-sum-sharp}
\end{equation}
which is the sharp Prouhet value.  The next degree is already explicit:
\begin{equation}
 S_{m,m+1}
 =(-1)^m2^{\binom m2}(m+1)!\frac{2^m-1}{2}.
 \label{eq:power-sum-next}
\end{equation}

\subsection{A Bell-polynomial compression}

Formula \eqref{eq:power-sum-composition} is finite and elementary, but its
number of summands grows.  A complete Bell polynomial compresses all
compositions into one standard object.  Let \(B_{2\ell}\) denote Bernoulli
numbers in the convention \(B_2=1/6\), and define
\begin{align}
 \kappa_1^{(m)}&=\frac{2^m-1}{2},
 \label{eq:kappa-one}\\
 \kappa_{2\ell}^{(m)}
 &=\frac{B_{2\ell}}{2\ell}
   \frac{2^{2\ell m}-1}{2^{2\ell}-1}
   \qquad(\ell\ge1),
 \label{eq:kappa-even}\\
 \kappa_{2\ell+1}^{(m)}&=0
   \qquad(\ell\ge1).
 \label{eq:kappa-odd}
\end{align}
Let \(\Bell_r\) be the complete exponential Bell polynomial, defined by
\begin{equation}
 \exp\!\left(\sum_{j\ge1}\kappa_j\frac{x^j}{j!}\right)
 =\sum_{r\ge0}\Bell_r(\kappa_1,\ldots,\kappa_r)\frac{x^r}{r!}.
 \label{eq:Bell-def}
\end{equation}

\begin{theorem}[Bell-polynomial power sums]
\label{thm:power-sum-bell}
For \(m\ge1\) and \(r\ge0\),
\begin{equation}
 \boxed{\quad
 S_{m,m+r}
 =(-1)^m2^{\binom m2}\frac{(m+r)!}{r!}
 \Bell_r\!\left(\kappa_1^{(m)},\ldots,\kappa_r^{(m)}\right).
 \quad}
 \label{eq:power-sum-bell}
\end{equation}
\end{theorem}

\begin{proof}
Use the classical formal expansion
\begin{equation}
 \log\!\left(\frac{\e^y-1}{y}\right)
 =\frac y2+\sum_{\ell\ge1}
   \frac{B_{2\ell}}{2\ell(2\ell)!}y^{2\ell}.
 \label{eq:log-expm1}
\end{equation}
Factoring \(-2^jx\) from each term in \eqref{eq:moment-egf} gives
\begin{equation}
 \prod_{j=0}^{m-1}(1-\e^{2^jx})
 =(-1)^m2^{\binom m2}x^m
   \exp\!\left(\sum_{r\ge1}\kappa_r^{(m)}\frac{x^r}{r!}\right).
 \label{eq:moment-bell-factor}
\end{equation}
Compare coefficients with \eqref{eq:Bell-def} and
\eqref{eq:moment-egf}.
\end{proof}

The first three surviving moments are therefore
\begin{align}
 S_{m,m}
 &=(-1)^m2^{\binom m2}m!,
 \label{eq:first-surviving}\\
 S_{m,m+1}
 &=(-1)^m2^{\binom m2}(m+1)!\kappa_1^{(m)},
 \label{eq:second-surviving}\\
 S_{m,m+2}
 &=(-1)^m2^{\binom m2}\frac{(m+2)!}{2}
   \left((\kappa_1^{(m)})^2+\kappa_2^{(m)}\right),
 \label{eq:third-surviving}
\end{align}
where \(\kappa_2^{(m)}=(4^m-1)/36\).

\begin{center}
\begin{tabular}{@{}c|rrrrrr@{}}
\toprule
\(m\backslash r\) & 0 & 1 & 2 & 3 & 4 & 5\\
\midrule
1 & 0 & -1 & -1 & -1 & -1 & -1\\
2 & 0 & 0 & 4 & 18 & 64 & 210\\
3 & 0 & 0 & 0 & -48 & -672 & -6720\\
4 & 0 & 0 & 0 & 0 & 1536 & 57600\\
5 & 0 & 0 & 0 & 0 & 0 & -122880\\
\bottomrule
\end{tabular}
\end{center}
The triangular zero region is Prouhet cancellation; the diagonal is
\eqref{eq:power-sum-sharp}.

\subsection{Binomial-basis moments}

Ordinary powers are not the only useful basis.  Define
\begin{equation}
 N_{m,r}=\sum_{n<2^m}\varepsilon_n\binom nr.
 \label{eq:newton-moment-def}
\end{equation}
Substituting \(z=1+u\) in \eqref{eq:finite-product} gives the exact coefficient
formula
\begin{equation}
 \boxed{\quad
 N_{m,r}=[u^r]\prod_{j=0}^{m-1}\left(1-(1+u)^{2^j}\right).
 \quad}
 \label{eq:newton-moment-product}
\end{equation}
Every factor is divisible by \(u\), so
\begin{equation}
 N_{m,r}=0\quad(r<m),
 \qquad
 N_{m,m}=(-1)^m2^{\binom m2}.
 \label{eq:newton-leading}
\end{equation}
The next coefficient is
\begin{equation}
 N_{m,m+1}
 =(-1)^m2^{\binom m2}\frac{2^m-1-m}{2}.
 \label{eq:newton-next}
\end{equation}
These identities are the Newton-series analogue of Prouhet extraction and can
be preferable in discrete proofs because finite differences diagonalize the
binomial basis.

\section{Substitution, automata, tensor products, and Walsh characters}
\label{sec:automatic}

\subsection{The morphic fixed point}

Let \(W_m\) be the zero--one word
\(\tau(0)\tau(1)\cdots\tau(2^m-1)\).  The recurrences
\eqref{eq:evenodd-bit} give
\begin{equation}
 W_{m+1}=W_m\,\overline{W_m},
 \label{eq:word-complement}
\end{equation}
where the bar complements every bit.  Equivalently, the infinite word is the
fixed point beginning in \(0\) of the uniform morphism
\begin{equation}
 0\longmapsto01,
 \qquad
 1\longmapsto10.
 \label{eq:morphism}
\end{equation}
This is the classical Thue construction \cite{thue1906,alloucheShallit1999}.
For signed words \(V_m=(\varepsilon_0,\ldots,\varepsilon_{2^m-1})\),
\begin{equation}
 V_{m+1}=(V_m,-V_m).
 \label{eq:signed-word-recursion}
\end{equation}

\subsection{A two-state matrix representation}

Let
\begin{equation}
 M_0=\begin{pmatrix}1&0\\0&1\end{pmatrix},
 \qquad
 M_1=\begin{pmatrix}0&1\\1&0\end{pmatrix},
 \qquad
 e_0=\binom10.
 \label{eq:automaton-matrices}
\end{equation}
For any binary expansion \(n=\sum_{j=0}^{m-1}b_j2^j\), define
\begin{equation}
 v(n)=M_{b_0}M_{b_1}\cdots M_{b_{m-1}}e_0.
 \label{eq:state-vector}
\end{equation}
The matrices commute, and each \(M_1\) toggles the two states.  Hence
\begin{equation}
 v(n)=\binom{1-\tau(n)}{\tau(n)},
 \qquad
 \tau(n)=\begin{pmatrix}0&1\end{pmatrix}v(n),
 \qquad
 \varepsilon_n=\begin{pmatrix}1&-1\end{pmatrix}v(n).
 \label{eq:matrix-output}
\end{equation}
This is the minimal deterministic automaton in linear-algebra form.

\subsection{Kronecker and Hadamard formulae}

Let \(u=(1,-1)^{\mathsf T}\).  In the natural binary order,
\begin{equation}
 \begin{pmatrix}
 \varepsilon_0\\ \varepsilon_1\\ \vdots\\ \varepsilon_{2^m-1}
 \end{pmatrix}
 =u^{\otimes m}.
 \label{eq:kronecker}
\end{equation}
Thus the length-\(2^m\) block is a row (or column, depending on convention) of
the Sylvester Hadamard matrix.  If \(\bm b(n)\in\Ftwo^m\) is the binary digit
vector and \(\bm 1=(1,\ldots,1)\), then
\begin{equation}
 \varepsilon_n=(-1)^{\bm1\cdot\bm b(n)}.
 \label{eq:walsh-character}
\end{equation}
This is the Walsh character indexed by the all-ones vector.

\begin{theorem}[Walsh transform is a point mass]
\label{thm:walsh-transform}
For \(\bm a\in\Ftwo^m\),
\begin{equation}
 \sum_{n=0}^{2^m-1}\varepsilon_n
 (-1)^{\bm a\cdot\bm b(n)}
 =\begin{cases}
 2^m,&\bm a=\bm1,\\
 0,&\bm a\ne\bm1.
 \end{cases}
 \label{eq:walsh-transform}
\end{equation}
\end{theorem}

\begin{proof}
The summand is the character indexed by \(\bm a+\bm1\).  Summing a nontrivial
character over the finite Boolean group gives zero, while the trivial
character sums to \(2^m\).
\end{proof}

\subsection{Rademacher-function form}

For \(x\in[0,1)\), define the right-continuous Rademacher functions
\begin{equation}
 r_j(x)=(-1)^{\lfloor2^{j+1}x\rfloor}
 \qquad(j\ge0).
 \label{eq:rademacher}
\end{equation}
If \(0\le n<2^m\) and \(x\in[n/2^m,(n+1)/2^m)\), the first \(m\) binary
digits of \(x\) are the \(m\)-bit expansion of \(n\).  Therefore
\begin{equation}
 \varepsilon_n=\prod_{j=0}^{m-1}r_j(x)
 \qquad\left(\frac{n}{2^m}\le x<\frac{n+1}{2^m}\right).
 \label{eq:rademacher-product}
\end{equation}
The signed Thue--Morse block is thus the step-function sampling of a Walsh
function.  This is the bridge from automatic sequences to harmonic analysis.

\section{Discrete Fourier formulae and Riesz products}
\label{sec:fourier}

\subsection{Exact finite Fourier transform}

For integer \(k\), define the length-\(2^m\) discrete Fourier transform
\begin{equation}
 \widehat\varepsilon_m(k)
 =\sum_{n=0}^{2^m-1}\varepsilon_n
   \exp\!\left(-\frac{2\pi\ii kn}{2^m}\right).
 \label{eq:dft-def}
\end{equation}
Evaluating \eqref{eq:finite-product} at a root of unity gives an exact
factorization.

\begin{theorem}[Finite Fourier product]
\label{thm:finite-fourier}
For every \(m\ge0\) and integer \(k\),
\begin{equation}
 \widehat\varepsilon_m(k)
 =\prod_{j=0}^{m-1}
  \left(1-\exp\!\left(-\frac{2\pi\ii k2^j}{2^m}\right)\right).
 \label{eq:dft-product}
\end{equation}
If \(m\ge1\), this vanishes for every even \(k\).  For odd \(k\),
\begin{equation}
 \widehat\varepsilon_m(k)
 =(2\ii)^m\exp\!\left(-\pi\ii k(1-2^{-m})\right)
 \prod_{r=1}^m\sin\!\left(\frac{\pi k}{2^r}\right),
 \label{eq:dft-sine}
\end{equation}
and hence
\begin{equation}
 \left|\widehat\varepsilon_m(k)\right|
 =2^m\prod_{r=1}^m
   \left|\sin\!\left(\frac{\pi k}{2^r}\right)\right|.
 \label{eq:dft-magnitude}
\end{equation}
\end{theorem}

\begin{proof}
Equation \eqref{eq:dft-product} is \eqref{eq:finite-product} with
\(z=\exp(-2\pi\ii k/2^m)\).  If \(k\) is even, one factor has exponent an
integer multiple of \(2\pi\ii\).  For odd \(k\), use
\(1-\e^{-\ii\theta}=2\ii\e^{-\ii\theta/2}\sin(\theta/2)\) in every factor
and sum the phases.
\end{proof}

The vanishing of every even Fourier mode is a spectral restatement of the
adjacent-pair cancellation.  The nonzero odd modes are controlled by a
lacunary sine product.

\subsection{Riesz-product densities}

On the unit circle, define the normalized squared amplitude
\begin{equation}
 R_m(x)=2^{-m}\left|P_m(\e^{2\pi\ii x})\right|^2.
 \label{eq:Riesz-def}
\end{equation}
Because \(|1-\e^{\ii\theta}|^2=2(1-\cos\theta)\),
\begin{equation}
 \boxed{\quad
 R_m(x)=\prod_{j=0}^{m-1}
       \left(1-\cos(2\pi2^jx)\right).
 \quad}
 \label{eq:Riesz-product}
\end{equation}
Parseval's identity gives
\begin{equation}
 \int_0^1R_m(x)\dd x=1,
 \label{eq:Riesz-mass}
\end{equation}
so \(R_m(x)\dd x\) is a probability measure.  Its weak limit is the classical
singular continuous Thue--Morse spectral measure; the exact finite formula
\eqref{eq:Riesz-product} is the part needed for algebraic and computational
work \cite{queffelec2010}.

\section{Infinite products with Thue--Morse exponents}
\label{sec:infinite-products}

The sequence can also be encoded as exponents in convergent products.  The
basic mechanism is again even--odd pairing.

\subsection{The pair-reduction operator}

Suppose \(R(n)\) is nonzero and the products converge.  Then
\begin{equation}
 \prod_{n\ge0}R(n)^{\varepsilon_n}
 =\prod_{q\ge0}
   \left(\frac{R(2q)}{R(2q+1)}\right)^{\varepsilon_q}.
 \label{eq:product-pair-reduction}
\end{equation}
A factor close only to \(1+O(1/n)\) is converted into a paired factor close
to \(1+O(1/n^2)\), which explains the favorable convergence of balanced
rational products.

Following Allouche, Riasat, and Shallit \cite{alloucheRiasatShallit2017}, let
\begin{equation}
 f(a,b)=\prod_{n\ge1}\left(\frac{n+a}{n+b}\right)^{\varepsilon_n},
 \label{eq:f-product}
\end{equation}
for parameters avoiding poles.  Pair reduction gives
\begin{equation}
 f(a,b)
 =\frac{1+b}{1+a}\,
   \frac{f(a/2,(a+1)/2)}{f(b/2,(b+1)/2)}.
 \label{eq:f-reduction}
\end{equation}
Defining
\begin{equation}
 g(x)=\frac{f(x/2,(x+1)/2)}{x+1},
 \label{eq:g-def}
\end{equation}
the same identity becomes
\begin{equation}
 f(a,b)=\frac{g(a)}{g(b)},
 \qquad
 (1+x)g(x)=\frac{g(x/2)}{g((x+1)/2)}.
 \label{eq:g-functional}
\end{equation}

Two striking specializations are
\begin{align}
 \prod_{n\ge0}
  \left(\frac{2n+1}{2n+2}\right)^{\varepsilon_n}
 &=\frac{1}{\sqrt2},
 \label{eq:woods-robbins}\\
 \prod_{n\ge0}
  \left(\frac{4n+1}{4n+3}\right)^{\varepsilon_n}
 &=\frac12.
 \label{eq:quarter-product}
\end{align}
The first is the Woods--Robbins product \cite{woods1978,robbins1979}.  These
identities are not individual-term formulae, but they are exact global
fingerprints of the same parity character and are generated by the same
binary recurrence.

\section{A base-\texorpdfstring{\(B\)}{B}, root-of-unity discovery engine}
\label{sec:generalization}

The product method is not peculiar to base two.  It is useful to state the
general pattern because it explains which Thue--Morse formulae are structural
and which are genuinely binary.

Let \(B\ge2\), let \(s_B(n)\) be the sum of the base-\(B\) digits of \(n\),
and fix \(\zeta\in\C\).  Define
\begin{equation}
 a_n=\zeta^{s_B(n)}.
 \label{eq:general-digit-character}
\end{equation}
Then
\begin{equation}
 a_{Bq+r}=\zeta^r a_q
 \qquad(0\le r<B),
 \label{eq:general-recurrence}
\end{equation}
and, for \(0\le n<B^m\), unique base-\(B\) expansion gives
\begin{equation}
 \sum_{n=0}^{B^m-1}a_nz^n
 =\prod_{j=0}^{m-1}
   \left(1+\zeta z^{B^j}+\zeta^2z^{2B^j}
      +\cdots+\zeta^{B-1}z^{(B-1)B^j}\right).
 \label{eq:general-product}
\end{equation}
The ordinary signed Thue--Morse sequence is \(B=2\), \(\zeta=-1\).

If \(\zeta\) is a nontrivial \(B\)th root of unity, every factor of
\eqref{eq:general-product} vanishes at \(z=1\).  Consequently the exponential
sum obtained by setting \(z=\e^x\) has a zero of order at least \(m\) at the
origin, and the generalized digit character annihilates powers of degree
below \(m\).  Thus the Prouhet phenomenon is the root-of-unity shadow of a
more general digit-product factorization.

\section{Computational comparison and possible Lean additions}
\label{sec:formalization}

Different formulae are efficient in different proof contexts.  The following
table summarizes their input data and natural use.

\begin{center}
\begin{tabularx}{\textwidth}{@{}>{\raggedright\arraybackslash}p{0.23\textwidth}
>{\raggedright\arraybackslash}p{0.25\textwidth}X@{}}
\toprule
Representation & Required primitive & Best use\\
\midrule
Binary parity & binary digits / \(\wt(n)\) & executable evaluation and induction\\
Block formula \eqref{eq:block-multiplicativity} & quotient and remainder by \(2^m\) & dyadic decomposition and Fabius sums\\
Factorial formula \eqref{eq:factorial-sign} & \(\vtwo(n!)\) & number-theoretic rewrites\\
Central binomial formula \eqref{eq:central-binomial-formula} & \(\vtwo\binom{2n}{n}\) & Kummer/carry arguments\\
Sparse Pascal formula \eqref{eq:sparse-pascal-exp} & \(O(\log n)\) binomial entries & Pascal-triangle connections\\
Coefficient product \eqref{eq:infinite-product-E} & formal power series & generating functions and Mahler equations\\
Prefix formula \eqref{eq:prefix-compact} & parity of \(N\) and one smaller term & exact discrepancy\\
Finite difference \eqref{eq:finite-difference-master} & shift operators & Prouhet cancellation for arbitrary functions\\
Fourier product \eqref{eq:dft-product} & roots of unity & spectral and exponential-sum analysis\\
Ruler update \eqref{eq:ruler-update-sign} & trailing-zero count & streaming generation\\
\bottomrule
\end{tabularx}
\end{center}

The current repository already contains named theorems for the binary
recurrences, block concatenation, dyadic power cancellation, sharp Prouhet
extraction, iterated prefix sums, and the binomial--logarithm evaluator
\cite{proveit}.  A compact complementary module could formalize the following
statements.

\begin{enumerate}
\item \texttt{thueMorseSign\_xor}: equation \eqref{eq:xor-law}.
\item \texttt{thueMorseSign\_add\_valuation}: equation
      \eqref{eq:addition-law}, using Legendre's formula for factorial
      valuations or a base-two Kummer theorem.
\item \texttt{thueMorseSign\_centralChoose}: equation
      \eqref{eq:central-binomial-formula}.
\item \texttt{thueMorseSign\_pascalPowTwo}: equation
      \eqref{eq:sparse-pascal-exp}, using the modulo-two Lucas theorem.
\item \texttt{thueMorsePrefix\_closed}: equation
      \eqref{eq:prefix-compact}.
\item \texttt{thueMorsePolynomial\_eq\_prod}: equation
      \eqref{eq:finite-product} in a polynomial or power-series ring.
\item \texttt{thueMorseSum\_eq\_mixedDifference}: equation
      \eqref{eq:finite-difference-master} for an additive target.
\item \texttt{thueMorsePowerSum\_composition}: equation
      \eqref{eq:power-sum-composition} over \(\Q\).
\end{enumerate}

The first six are algebraic and finite.  They should require no analytic
infrastructure.  The Bell-polynomial compression is less urgent: the finite
composition formula is simpler to formalize and already determines every
coefficient.  The Fourier product is also finite and can be stated over
\(\C\) after the polynomial factorization is available.

\begin{boxedremark}
\textbf{Why these formulae matter for the Fabius development.}
The global Fabius extension uses \(\varepsilon_n\) as translate signs, while
the dyadic closed forms use complete blocks of those signs.  Block
multiplicativity shortens reindexing arguments; the carry law controls signs
when indices are added rather than concatenated; the finite-difference form
packages every polynomial cancellation into one operator identity; and the
all-order power-sum formula supplies exact coefficients beyond the first
surviving Prouhet moment.  None changes the Fabius function, but each offers a
new proof interface to the same arithmetic layer.
\end{boxedremark}

\section{Worked examples}
\label{sec:examples}

\subsection{One index through six formulae}

Take \(n=13=(1101)_2\).  Then \(\wt(13)=3\), so
\(\tau(13)=1\) and \(\varepsilon_{13}=-1\).  The alternatives give:
\begin{align*}
 \vtwo(13!)&=\left\lfloor\frac{13}{2}\right\rfloor
 +\left\lfloor\frac{13}{4}\right\rfloor
 +\left\lfloor\frac{13}{8}\right\rfloor=6+3+1=10,\\
 (-1)^{13-\vtwo(13!)}&=(-1)^3=-1,\\
 \vtwo\binom{26}{13}&=3,\\
 \binom{13}{1}+\binom{13}{2}+\binom{13}{4}+\binom{13}{8}
 &=13+78+715+1287=2093\equiv1\pmod2,\\
 \nu(13)&=2^3=8,\\
 [z^{13}]\prod_{j\ge0}(1-z^{2^j})&=-1.
\end{align*}
The sparse Pascal entries have parities \(1,0,1,1\), reproducing the bits of
\(13\) from least to most significant position.

\subsection{An all-order moment}

For \(m=3\) and \(r=5\), direct summation gives
\begin{equation*}
 S_{3,5}=\sum_{n=0}^7\varepsilon_nn^5=-6720.
\end{equation*}
Formula \eqref{eq:power-sum-composition} sums the positive compositions of
\(5\) into three parts:
\begin{equation*}
 (1,1,3),(1,2,2),(1,3,1),(2,1,2),(2,2,1),(3,1,1).
\end{equation*}
Thus
\begin{align*}
 S_{3,5}
 &=-5!\sum_{\substack{q_0+q_1+q_2=5\\ q_j\ge1}}
   \frac{2^{q_1+2q_2}}{q_0!q_1!q_2!}\\
 &=-6720.
\end{align*}
The formula is exact and uses only six rational summands; no recursion through
lower moments is required.

\subsection{A Fourier sample}

For \(m=4\), all even DFT frequencies vanish.  At \(k=1\),
\begin{equation*}
 \widehat\varepsilon_4(1)
 =(2\ii)^4\e^{-15\pi\ii/16}
 \sin\frac\pi2\sin\frac\pi4\sin\frac\pi8\sin\frac\pi{16}.
\end{equation*}
The product factorization exposes both the exact phase and the hierarchy of
dyadic scales.  Directly summing sixteen complex terms obscures that
structure.

\section{Conclusion}
\label{sec:conclusion}

The Thue--Morse term is not merely ``the parity of the binary digit sum.''  It
is simultaneously:
\begin{itemize}
\item a character of the Boolean group under exclusive-or;
\item a digit-block multiplicative sequence with an exact carry correction
      under ordinary addition;
\item the parity of a factorial valuation and of a central-binomial valuation;
\item the parity encoded by the power-of-two columns of Pascal's triangle;
\item the coefficient sequence of a lacunary product satisfying a Mahler
      equation;
\item a mixed finite-difference kernel whose dyadic moments admit closed
      all-order formulae;
\item a Walsh--Hadamard row and a Rademacher product;
\item a finite Fourier product and the seed of a Riesz product;
\item and the exponent pattern behind a family of self-similar infinite
      products.
\end{itemize}

The binomial--logarithm formula in the existing Fabius document is therefore
one member of a broad equivalence class.  Its distinctive feature is that it
recovers binary parity from the aggregate oddness pattern of an entire Pascal
row.  The sparse Pascal formula recovers the same parity from logarithmically
many distinguished entries; Kummer's formula recovers it from carries;
Legendre's formula recovers it from the twos in a factorial; and the Mahler
product recovers it as a coefficient.  The finite-difference identity then
explains why the same signs are so effective in exact Fabius calculations:
they are an inclusion--exclusion kernel adapted perfectly to the dyadic
steps \(1,2,4,\ldots\).

\appendix

\section{The first 64 terms}
\label{app:first-terms}

\begin{center}
\small
\begin{longtable}{@{}r|*{8}{c}@{}}
\toprule
\(n\) & \multicolumn{8}{c}{successive values}\\
\midrule
\endhead
0--7   & 0&1&1&0&1&0&0&1\\
8--15  & 1&0&0&1&0&1&1&0\\
16--23 & 1&0&0&1&0&1&1&0\\
24--31 & 0&1&1&0&1&0&0&1\\
32--39 & 1&0&0&1&0&1&1&0\\
40--47 & 0&1&1&0&1&0&0&1\\
48--55 & 0&1&1&0&1&0&0&1\\
56--63 & 1&0&0&1&0&1&1&0\\
\bottomrule
\end{longtable}
\end{center}
These rows display \(\tau(n)\).  Replacing every \(0\) by \(+1\) and every
\(1\) by \(-1\) gives \(\varepsilon_n\).  Each row of length eight is a copy
or complement according to the sign of its block index, as predicted by
\eqref{eq:block-multiplicativity}.

\section{A compact identity map}
\label{app:identity-map}

For quick reference, the main transformations are
\begin{align*}
 \tau(n)
 &\longleftrightarrow \varepsilon_n=1-2\tau(n),\\
 \varepsilon_n
 &\longleftrightarrow (-1)^{\wt(n)},\\
 \wt(n)
 &\longleftrightarrow n-\vtwo(n!),\\
 \wt(n)
 &\longleftrightarrow \vtwo\binom{2n}{n},\\
 b_j(n)
 &\longleftrightarrow \binom{n}{2^j}\bmod2,\\
 \varepsilon_n
 &\longleftrightarrow [z^n]\prod_{j\ge0}(1-z^{2^j}),\\
 \sum_{n<2^m}\varepsilon_nf(x+n)
 &\longleftrightarrow (-1)^m\Delta_1\Delta_2\cdots
   \Delta_{2^{m-1}}f(x),\\
 P_m(\e^{-2\pi\ii k/2^m})
 &\longleftrightarrow \widehat\varepsilon_m(k).
\end{align*}
Each arrow is an exact identity, not an asymptotic or heuristic analogy.

\begin{thebibliography}{99}

\bibitem{prouhet1851}
E.~Prouhet,
\newblock M\'emoire sur quelques relations entre les puissances des nombres,
\newblock \emph{Comptes rendus de l'Acad\'emie des Sciences de Paris}
  \textbf{33} (1851), 225.

\bibitem{thue1906}
A.~Thue,
\newblock \emph{\"Uber unendliche Zeichenreihen},
\newblock Norske Videnskabers Selskabs Skrifter, I. Matematisk-naturvidenskabelig
  Klasse, no.~7 (1906), 1--22.

\bibitem{morse1921}
M.~Morse,
\newblock Recurrent geodesics on a surface of negative curvature,
\newblock \emph{Transactions of the American Mathematical Society}
  \textbf{22} (1921), 84--100.

\bibitem{kummer1852}
E.~E.~Kummer,
\newblock \"Uber die Erg\"anzungss\"atze zu den allgemeinen
  Reciprocit\"atsgesetzen,
\newblock \emph{Journal f\"ur die reine und angewandte Mathematik}
  \textbf{44} (1852), 93--146.

\bibitem{lucasSurvey}
R.~Me\v{s}trovi\'c,
\newblock Lucas' theorem: its generalizations, extensions and applications
  (1878--2014),
\newblock arXiv:\nolinkurl{1409.3820} (2014).

\bibitem{fine1947}
N.~J.~Fine,
\newblock Binomial coefficients modulo a prime,
\newblock \emph{American Mathematical Monthly} \textbf{54} (1947), 589--592.

\bibitem{mahler1929}
K.~Mahler,
\newblock Arithmetische Eigenschaften der L\"osungen einer Klasse von
  Funktionalgleichungen,
\newblock \emph{Mathematische Annalen} \textbf{101} (1929), 342--366.

\bibitem{woods1978}
D.~R.~Woods,
\newblock Elementary problem E2692,
\newblock \emph{American Mathematical Monthly} \textbf{85} (1978), 48.

\bibitem{robbins1979}
D.~Robbins,
\newblock Solution to elementary problem E2692,
\newblock \emph{American Mathematical Monthly} \textbf{86} (1979), 394--395.

\bibitem{alloucheShallit1999}
J.-P.~Allouche and J.~Shallit,
\newblock The ubiquitous Prouhet--Thue--Morse sequence,
\newblock in C.~Ding, T.~Helleseth, and H.~Niederreiter (eds.),
  \emph{Sequences and Their Applications}, Springer, 1999, 1--16.

\bibitem{allouche2003}
J.-P.~Allouche and J.~Shallit,
\newblock \emph{Automatic Sequences: Theory, Applications, Generalizations},
\newblock Cambridge University Press, 2003.

\bibitem{queffelec2010}
M.~Queff\'elec,
\newblock \emph{Substitution Dynamical Systems -- Spectral Analysis},
\newblock second edition, Lecture Notes in Mathematics~1294, Springer, 2010.

\bibitem{alloucheRiasatShallit2017}
J.-P.~Allouche, S.~Riasat, and J.~Shallit,
\newblock More infinite products: Thue--Morse and the Gamma function,
\newblock arXiv:\nolinkurl{1709.03398v2} (2017).

\bibitem{proveitSynthesis}
V.~Reshetnikov,
\newblock \emph{The Fabius Function and Rvachev's Up-Function: Exact
  arithmetic, probability, Fourier analysis, and corrected sharp
  asymptotics},
\newblock source in \path{Analysis/FabiusFunction/docs/Fabius_Function_and_Rvachev_Up/},
  consulted 25 August 2026.

\bibitem{proveit}
V.~Reshetnikov,
\newblock \emph{ProveIt: Analysis/FabiusFunction},
\newblock Lean~4 formal development, consulted 25 August 2026.
\newblock \url{https://github.com/VladimirReshetnikov/ProveIt/tree/main/Analysis/FabiusFunction}.

\end{thebibliography}

\end{document}
